\documentclass[11pt]{article}
\usepackage[a4paper,margin=28mm]{geometry}
\usepackage{amsmath,amssymb,amsthm,booktabs,array,longtable}
\usepackage[T1]{fontenc}
\usepackage[colorlinks=true,linkcolor=blue,urlcolor=blue,citecolor=blue]{hyperref}
\hypersetup{pdftitle={Mathar's half-factorial recurrence conjectures: extensions of Jovovic's formulas and Barry's array sums},pdfauthor={Haobo Ma and Wenlin Zhang}}
\newtheorem{theorem}{Theorem}
\newtheorem{lemma}{Lemma}
\newtheorem{corollary}{Corollary}
\DeclareMathOperator{\erf}{erf}
\title{Mathar's half-factorial recurrence conjectures:\\
extensions of Jovovic's formulas and Barry's array sums}
\author{Haobo Ma\\
  \small ChronoAI Pte Ltd, Singapore
  \and Wenlin Zhang\\
  \small National University of Singapore\\
  \small ChronoAI Pte Ltd, Singapore}
\date{15 November 2026}
\begin{document}
\maketitle
\begin{abstract}
Extending Jovovic's generating-function formulas of 25 September 2003
and Barry's array-sum identity (2009, Section~13.2, pp.~283--284),
we prove Mathar's recurrence conjectures in the Formula fields of
OEIS A081124, revision~8 (24 November 2012), and A084261, revision~9
(30 November 2012), retained in revisions~20 and~29, respectively.
A common finite-sum argument proves both recurrences
for every $n\geq3$. For the fixed-point polynomials of Fang, Hamaker and
Troyka, we also give a parameter-dependent recurrence and self-contained
proofs of the continued fractions and Hankel determinants recorded by
Hanna and Barry. Further consequences include partial sums and linked
factorial-ratio identities.
\end{abstract}

\section{The two recurrences}
For $n\geq0$, put
\[
 a_n=\sum_{k=0}^{n}\binom nk\lfloor k/2\rfloor!,\qquad
 c_n=\sum_{j=0}^{\lfloor n/2\rfloor}\binom n{2j}j!.
\]
Mathar conjectured the recurrences below in OEIS
A081124 and its directly linked entry A084261, respectively
\cite{oeis-a,oeis-c}. Their leading coefficients are nonzero, so each
recurrence and its three initial values determine the entire sequence.
Our proof treats the two sequences together, by separating the parity
of the factorial index before taking the binomial transform.

The binomial-sum descriptions already have combinatorial interpretations.
Goyt and Pudwell's Theorem~14 identifies $a_n$ with a class of
two-colored pattern-avoiding set partitions \cite[Section~4.5, p.~13]{gp}.
Fang, Hamaker and Troyka obtain the sum defining $c_n$ in the example
following their Proposition~4.13 \cite[Section~4.3, p.~16]{fht}.
These earlier proofs identify the counting sequences and, in the second
case, the polynomials marking fixed points. Barry's thesis already
identifies the array sums giving $a_n$
\cite[Section~13.2, pp.~283--284]{barry-thesis} and the $t=2$ instance
of the continued fraction below \cite[Example~64, pp.~53--54]{barry-thesis}.
We extend these formulas by establishing Mathar's recurrences and a
weighted recurrence, with full proofs of the related identities.
The generating functions in Section~\ref{sec:egf} are Jovovic's formulas;
the continued fractions and Hankel products are likewise attributed at
their statements.

Equivalent second-order and partial-sum formulas for $a_n$ were recorded
by Sequence Machine in 2021 \cite{sequence-machine}; both four-term
recurrences occur in Nuspl's 2023 recurrence data \cite{nuspl}.
These records concern formulas agreeing with sequence terms.
The proofs below establish the identities from the defining sums.

\begin{theorem}\label{thm:recurrences}
For every integer $n\geq3$,
\begin{align}
 2a_n-4a_{n-1}+(2-n)a_{n-2}+(n-1)a_{n-3}&=0,
 \label{eq:a}\\
 2c_n-4c_{n-1}+(2-n)c_{n-2}+(n-2)c_{n-3}&=0.
 \label{eq:c}
\end{align}
The initial values are $(a_0,a_1,a_2)=(1,2,4)$ and
$(c_0,c_1,c_2)=(1,1,2)$.
\end{theorem}

\section{A common finite-sum proof}
\begin{lemma}\label{lem:transfer}
Let $q\in\mathbb Q$ and $u_0,u_1,\ldots\in\mathbb Q$ satisfy
\begin{equation}\label{eq:input}
 2u_{k+3}+2u_{k+2}-(k+3)u_{k+1}-(k+q)u_k=0
 \qquad(k\geq0).
\end{equation}
Then $v_n=\sum_{k=0}^{n}\binom nk u_k$ satisfies
\begin{equation}\label{eq:transfer}
 2v_{n+3}-4v_{n+2}-(n+1)v_{n+1}+(n+3-q)v_n=0
 \qquad(n\geq0).
\end{equation}
\end{lemma}
\begin{proof}
Define $U_r(n)=\sum_{k=0}^{n}\binom nk u_{k+r}$ for $r\geq0$.
Pascal's identity, with out-of-range binomial coefficients set to zero,
gives
\[
 U_r(n+1)=U_r(n)+U_{r+1}(n).
\]
Consequently, for fixed $n$, writing $v_{n+j}$ without its argument,
\begin{align*}
 U_0(n)&=v_n, & U_1(n)&=v_{n+1}-v_n,\\
 U_2(n)&=v_{n+2}-2v_{n+1}+v_n,&
 U_3(n)&=v_{n+3}-3v_{n+2}+3v_{n+1}-v_n.
\end{align*}
The weighted terms in \eqref{eq:input} obey
\begin{equation}\label{eq:weighted}
 \sum_{k=0}^{n}\binom nk k(u_{k+1}+u_k)=nU_1(n).
\end{equation}
Indeed, $(n-k)\binom nk=(k+1)\binom n{k+1}$, whence
\begin{align*}
 nU_1(n)-\sum_{k=0}^{n}\binom nk k u_{k+1}
 &=\sum_{k=0}^{n}\binom n{k+1}(k+1)u_{k+1}\\
 &=\sum_{k=0}^{n}\binom nk k u_k.
\end{align*}
The last reindexing adds the zero term at $k=0$ and removes the zero
term with binomial coefficient $\binom n{n+1}$; it also applies when
$n=0$. Multiply \eqref{eq:input} by $\binom nk$ and sum. By
\eqref{eq:weighted}, the resulting identity is
\[
 2U_3(n)+2U_2(n)-3U_1(n)-qU_0(n)=nU_1(n).
\]
Substitution of the four displayed expressions for $U_r(n)$ gives
\eqref{eq:transfer}.
\end{proof}

\begin{proof}[Proof of Theorem~\ref{thm:recurrences}]
For the first sequence take $u_{2m}=u_{2m+1}=m!$ and $q=1$.
For $k=2m$ the left-hand side of \eqref{eq:input} is
\[
 4(m+1)!-(4m+4)m!=0.
\]
For $k=2m+1$ it is
\[
 2(m+2)!-(2m+2)(m+1)!-(2m+2)m!
 =2(m+1)!-(2m+2)m!=0.
\]
Applying Lemma~\ref{lem:transfer} with $q=1$ at index $n-3$ gives
\eqref{eq:a} for every $n\geq3$.

For the second sequence take $u_{2m}=m!$, $u_{2m+1}=0$, and $q=2$.
The two parity cases of \eqref{eq:input} reduce to
\[
 2(m+1)!-(2m+2)m!=0,\qquad
 2(m+2)!-(2m+4)(m+1)!=0,
\]
respectively. Applying Lemma~\ref{lem:transfer} with $q=2$ at index
$n-3$ gives \eqref{eq:c} for every $n\geq3$.
Direct evaluation at $n=0,1,2$ gives the stated initial values.
\end{proof}

\section{Generating functions and partial sums}
\label{sec:egf}
\begin{corollary}\label{cor:egf}
Let $H(x)=\sqrt\pi\,e^{x^2/4}\erf(x/2)$, where
$\erf(z)=\frac{2}{\sqrt\pi}\int_0^z e^{-t^2}\,dt$.
The exponential generating functions are
\begin{align}
 A(x):=\sum_{n\geq0}a_n\frac{x^n}{n!}
 &=e^x\bigl(1+(1+x/2)H(x)\bigr),\label{eq:egfa}\\
 C(x):=\sum_{n\geq0}c_n\frac{x^n}{n!}
 &=e^x\bigl(1+(x/2)H(x)\bigr).\label{eq:egfc}
\end{align}
They are the formulas recorded by Jovovic on 25 September 2003
in A081124 and A084261. Moreover, for every $n\geq0$,
\begin{equation}\label{eq:relation}
 c_{n+1}=a_{n+1}-a_n,\qquad a_n=\sum_{j=0}^{n}c_j.
\end{equation}
\end{corollary}
\begin{proof}
Consider the odd entire series
$H_0(x)=\sum_{m\geq0}m!x^{2m+1}/(2m+1)!$.
The ratio of successive nonzero terms tends to zero for every fixed
$x$, so termwise differentiation is valid. Since
$m!/(2m)!=(m-1)!/(2(2m-1)!)$ for $m\geq1$, we have
\[
 \sum_{m\geq0}\frac{m!x^{2m}}{(2m)!}=1+\frac x2H_0(x),
 \qquad H_0'(x)=1+\frac x2H_0(x),\quad H_0(0)=0.
\]
Multiplication by $e^{-x^2/4}$ and integration from $0$ to $x$
give $H_0(x)=e^{x^2/4}\int_0^x e^{-t^2/4}\,dt=H(x)$.
For any sequence $u$, multiplication of its exponential generating
function by $e^x$ takes its coefficient $u_k$ to
$\sum_{k=0}^{n}\binom nk u_k$, as follows by taking the coefficient
of $x^n/n!$ in the Cauchy product. Applying this to the even series
$1+xH(x)/2$ and to its sum with $H(x)$ proves
\eqref{eq:egfa}--\eqref{eq:egfc}.

Put $D=A-C=e^xH$. Using $H'=1+xH/2$ gives
\[
 D'=e^x\bigl(H+H'\bigr)
    =e^x\bigl(1+(1+x/2)H\bigr)=A.
\]
Taking coefficients of $x^n/n!$ gives $a_{n+1}-c_{n+1}=a_n$.
Since $a_0=c_0=1$, summation gives the second identity in
\eqref{eq:relation}.
\end{proof}

\begin{corollary}\label{cor:shorter}
For every $n\geq2$,
\begin{equation}\label{eq:shorter}
 c_n=1+\frac n2a_{n-2},\qquad
 a_n=1+a_{n-1}+\frac n2a_{n-2}.
\end{equation}
Let $B(i,j)=\sum_{k\geq0}\binom ik\binom jk k!$ for integers $i,j\geq0$,
where out-of-range binomial coefficients vanish. Then
\begin{equation}\label{eq:somos}
 a_n=\sum_{i+j=n}B(i,j)\qquad(n\geq0).
\end{equation}
This is Somos's symmetric array A088699 \cite{somos}; its antidiagonal
sum identity was also established by Barry
\cite[Section~13.2, pp.~283--284]{barry-thesis}.
\end{corollary}
\begin{proof}
The identity $D'=A$, $D(0)=0$ from Corollary~\ref{cor:egf} and
$C=e^x+xD/2$ give $c_n=1+(n/2)a_{n-2}$ for $n\geq2$.
Using $c_n=a_n-a_{n-1}$ gives the second equality.
It is the second-order formula recorded by Sequence Machine on
9 December 2021; the partial-sum formula \eqref{eq:relation} is recorded
there on 1 November 2021 \cite{sequence-machine}.
For the array, Vandermonde's identity in its summed form gives
\[
 \sum_{i+j=n}B(i,j)
 =\sum_{k\geq0}k!\binom{n+1}{2k+1}
 =\sum_{j=0}^nc_j=a_n.
\]
The first equality uses
$\sum_{i=0}^n\binom ik\binom{n-i}k=\binom{n+1}{2k+1}$:
count a $(2k+1)$-element subset by its middle element $i+1$.
The second equality uses
$\sum_{j=2k}^n\binom j{2k}=\binom{n+1}{2k+1}$.

There is also a direct array explanation of the shorter recurrence.
$B(i,j)$ counts partial matchings between sets of sizes $i$ and $j$.
An element of the first set is either unmatched or paired to one of
the $j$ elements of the second set, so
$B(i,j)=B(i-1,j)+jB(i-1,j-1)$ for $i>0$, with $B(0,j)=1$.
Summing on $i+j=n$ leaves the boundary term $1$ and the previous
antidiagonal sum. The remaining term is
$\sum_{h+k=n-2}(k+1)B(h,k)$. Since $B(h,k)=B(k,h)$, it equals
$(n/2)\sum_{h+k=n-2}B(h,k)=(n/2)a_{n-2}$.
This also proves the second equality of \eqref{eq:shorter}.
Subtracting its instances at $n$ and $n-1$ recovers \eqref{eq:a}.
\end{proof}

\section{Weighted transforms, continued fractions and Hankel determinants}
For an indeterminate $t$, define
\[
 c_n(t)=\sum_{j=0}^{\lfloor n/2\rfloor}\binom n{2j}j!t^{n-2j},\qquad
 a_n(t)=\sum_{k=0}^{n}\binom nk\lfloor k/2\rfloor!t^{n-k}.
\]
Thus $c_n(1)=c_n$ and $a_n(1)=a_n$.
Fang, Hamaker and Troyka's fixed-point enumerator specializes to
$c_n(t)$ \cite[Proposition~4.13 and its following example]{fht}.
Its coefficient of $t^k$ is the triangle A161556:
\[
 T(n,k)=
 \begin{cases}\binom nk((n-k)/2)!,&0\leq k\leq n,\ n-k\text{ even},\\
 0,&\text{otherwise}.
 \end{cases}
\]
In particular, the row sums are $c_n$, as recorded by Barry
\cite{oeis-triangle}.

\begin{corollary}\label{cor:weighted}
For $n\geq3$,
\begin{equation}\label{eq:weighted-recurrence}
 2c_n(t)-4t c_{n-1}(t)+(2t^2-n)c_{n-2}(t)
       +(n-2)t c_{n-3}(t)=0,
\end{equation}
with $c_0(t)=1$, $c_1(t)=t$, $c_2(t)=t^2+1$. Also, for $n\geq0$,
\begin{equation}\label{eq:weighted-sum}
 a_n(t)=c_n(t)+\sum_{j=0}^{n-1}t^{n-1-j}c_j(t),
\end{equation}
where an empty sum is zero.
\end{corollary}
\begin{proof}
Put $E(x)=1+xH(x)/2$. The identity $H'=1+xH/2$ gives
$2E''-xE'-2E=0$. Since
$C_t(x):=\sum c_n(t)x^n/n!=e^{tx}E(x)$, it follows that
\[
 2C_t''-(x+4t)C_t'+(tx+2t^2-2)C_t=0.
\]
Comparing coefficients of $x^{n-2}/(n-2)!$ gives
\eqref{eq:weighted-recurrence}; the initial values follow from the sum.
Likewise $A_t-C_t=e^{tx}H$, where $A_t=\sum a_n(t)x^n/n!$.
Writing $D_t=A_t-C_t$ gives $D_t'=tD_t+C_t$ and $D_t(0)=0$.
The coefficient recurrence $d_{n+1}=td_n+c_n(t)$, with
$D_t=\sum d_nx^n/n!$, yields \eqref{eq:weighted-sum}.
The formula for $T(n,k)$ follows by setting $k=n-2j$ in $c_n(t)$.
\end{proof}

For a sequence $s_0,s_1,\ldots$, its Hankel determinant at index $n$ is
$\det(s_{i+j})_{0\leq i,j\leq n}$.
The following identities include Hanna's 17 January 2006 continued
fraction in A084261, Barry's continued fraction for A161556, and the
Hankel product A137704 \cite{oeis-c,oeis-triangle,oeis-hankel}.
All ordinary generating functions here are formal power series;
analytic convergence is not needed.

\begin{corollary}\label{cor:hankel}
Set $\lambda_j=\lfloor(j+1)/2\rfloor$ for $j\geq1$. Then
\begin{equation}\label{eq:jacobi}
 \sum_{n\geq0}c_n(t)z^n
 =\cfrac{1}{1-tz-\cfrac{\lambda_1z^2}{1-tz-
       \cfrac{\lambda_2z^2}{1-tz-\cfrac{\lambda_3z^2}{\ddots}}}},
\end{equation}
where $(\lambda_1,\lambda_2,\lambda_3,\ldots)=(1,1,2,2,3,3,\ldots)$.
For $n\geq0$ the Hankel determinants are independent of $t$ and equal
\begin{equation}\label{eq:hankel}
 \det(c_{i+j}(t))_{0\leq i,j\leq n}
 =\prod_{k=0}^{n}\left\lfloor\frac{k+2}{2}\right\rfloor^{n-k}.
\end{equation}
\end{corollary}
\begin{proof}
Let $u_{2m}=m!$, $u_{2m+1}=0$. The lower triangular matrix
$P_t=(\binom ir t^{i-r})_{0\leq i,r\leq n}$, with entries zero for $r>i$,
has determinant one. Vandermonde's identity gives
\[
 (c_{i+j}(t))_{i,j}=P_t(u_{i+j})_{i,j}P_t^{\mathsf T}.
\]
Indeed, the coefficient of $u_k t^{i+j-k}$ on the right is
$\sum_{r+s=k}\binom ir\binom js=\binom{i+j}k$.
Thus the two Hankel determinants agree. Reordering the even and odd
indices gives diagonal blocks $(i+j)!$ of size
$r=\lceil(n+1)/2\rceil$ and $(i+j+1)!$ of size
$s=\lfloor(n+1)/2\rfloor$. For $\alpha=0,1$,
\[
 \det((i+j+\alpha)!)_{0\leq i,j<m}
 =\prod_{j=0}^{m-1}j!(j+\alpha)!.
\]
To see this, factor $(i+\alpha)!$ from row $i$; the remaining column $j$
is the monic degree-$j$ polynomial
$(i+\alpha+1)\cdots(i+\alpha+j)$ evaluated at $i$.
Its determinant is the Vandermonde product
$\prod_{0\leq i<k<m}(k-i)=\prod_{j=0}^{m-1}j!$.
Empty products are one. Hence the Hankel determinant is
\[
 \Delta_n=\prod_{j=0}^{r-1}(j!)^2\prod_{j=0}^{s-1}j!(j+1)!.
\]
With $\Delta_{-1}=1$, its successive quotients are
$\Delta_{2m}/\Delta_{2m-1}=(m!)^2$ and
$\Delta_{2m+1}/\Delta_{2m}=m!(m+1)!$. Equivalently,
$\Delta_n\Delta_{n-2}/\Delta_{n-1}^2=\lfloor(n+1)/2\rfloor$
for $n\geq1$. Together with $\Delta_0=1$ this proves \eqref{eq:hankel}.

The same determinants give the continued fraction. The positive functional
\[
 L(f)=\int_{-\infty}^{\infty}f(y)|y|e^{-y^2}\,dy
\]
satisfies $L(y^{2m})=m!$ and $L(y^{2m+1})=0$; substitute $v=y^2$
and integrate by parts $m$ times for the first equality.
Since $L(p^2)>0$ for every nonzero real polynomial $p$,
Gram--Schmidt gives monic orthogonal polynomials $p_j$, with $p_0=1$.
Writing $h_j=L(p_j^2)$, triangular change of basis in the Gram matrix
gives $h_j=\Delta_j/\Delta_{j-1}$. Symmetry gives
$p_j(-y)=(-1)^jp_j(y)$. Moving multiplication by $y$ across the inner
product shows that all coefficients below $p_{j-1}$ and the coefficient
of $p_j$ in $yp_j$ vanish. The coefficient of $p_{j-1}$ is
$h_j/h_{j-1}=\lambda_j$. Thus
\[
 yp_0=p_1,\qquad yp_j=p_{j+1}+\lambda_jp_{j-1}\quad(j\geq1).
\]
For real $t$, $L((y+t)^n)=c_n(t)$. Expanding multiplication by $y+t$
in this basis expresses this moment as the total weight of paths from
height zero back to zero, staying nonnegative, with upward, level and
downward steps of respective weights $1$, $t$ and $\lambda_j$ (a down
step starts at height $j$). If $G_j(z)$ counts such paths above height
$j$, first-return decomposition gives
\[
 G_j=\frac{1}{1-tz-\lambda_{j+1}z^2G_{j+1}}.
\]
This proves \eqref{eq:jacobi}. Every coefficient involves finitely many
heights; polynomial dependence on $t$ extends the identities from real
$t$ to the indeterminate $t$.
\end{proof}

\section{Elementary identities in the linked entries}
The sequence $b_n=\lfloor n/2\rfloor!$ is A081123; its index
$\lfloor n/2\rfloor$ is A004526 \cite{oeis-b}. For $n\geq2$, the
two-part partitions of $n$ are $(j,n-j)$ with $1\leq j\leq\lfloor n/2\rfloor$.
Their smaller parts have product $b_n$, proving Hurt's comment in
A081123 (25 January 2013).

\begin{corollary}\label{cor:ratio}
Let $r_n=n!/\lfloor n/2\rfloor!$ (A081125). Its exponential generating
function is $(1+x)e^{x^2}$, as recorded by Jovovic on 24 September 2003.
Mathar's recorded recurrence of 26 November 2012 is valid for $n\geq3$:
\begin{equation}\label{eq:ratio}
 r_n+2r_{n-1}-2nr_{n-2}+4(2-n)r_{n-3}=0.
\end{equation}
For $n\geq1$, $r_n$ is the product of the larger parts in the two-part
partitions of $n+1$. Write
$\Gamma(z)=\int_0^\infty v^{z-1}e^{-v}\,dv$ for real $z>0$.
Its factorial, product, gamma and swinging-factorial forms are
\begin{align}
 r_n&=\prod_{i=\lfloor n/2\rfloor+1}^{n}i
 =\frac{2^n\Gamma((n+1+(n\bmod2))/2)}{\sqrt\pi},\label{eq:ratio-product}\\
 r_n&=b_n s_n=\sqrt{n!s_n},\qquad
 s_n:=\frac{n!}{(\lfloor n/2\rfloor!)^2}.\label{eq:swing}
\end{align}
Here $s_n$ is the swinging factorial A056040.
Eldar's two reciprocal sums of 10 March 2022 are
\begin{equation}\label{eq:reciprocal}
 \sum_{n\geq0}\frac1{r_n}=1+\frac32 H(1),\qquad
 \sum_{n\geq0}\frac{(-1)^n}{r_n}=1-\frac12 H(1).
\end{equation}
\end{corollary}
\begin{proof}
The even and odd terms give
$\sum r_nx^n/n!=\sum_{m\geq0}(x^{2m}+x^{2m+1})/m!=(1+x)e^{x^2}$.
The relations $r_{2m}=2r_{2m-1}$ for $m\geq1$ and
$r_{2m+1}=(2m+1)r_{2m}$ for $m\geq0$ prove \eqref{eq:ratio}:
for even $n$ its left side is
$2(n-2)(r_{n-2}-2r_{n-3})$; for odd $n$ it is
$4(r_{n-2}-(n-2)r_{n-3})$.
The larger-part product is
$n(n-1)\cdots(n+1-\lfloor(n+1)/2\rfloor)=r_n$.
The first equality of \eqref{eq:ratio-product} is factorial cancellation;
the gamma equality follows in each parity from
$\Gamma(m+1/2)=(2m)!\sqrt\pi/(4^mm!)$.
This half-integer formula follows by integration by parts from
$\Gamma(1/2)=\sqrt\pi$ (the Gaussian integral).
Equation~\eqref{eq:swing} follows from the definition of $s_n$ and $r_n>0$.
The same factorial quotient counts ordered selections of
$\lceil n/2\rceil$ distinct objects from $n$ objects.
Finally, $1/r_n=b_n/n!$, whose entire generating series is
$1+(1+x/2)H(x)$. Substitution of $x=1,-1$ proves
\eqref{eq:reciprocal}, including absolute convergence.
These are the formulas and comments of A081125 \cite{oeis-ratio}.
\end{proof}

One recorded interval product needs an endpoint correction. In A081125,
revision~86, the Formula block signed W.~I.~Hurt, 6 June 2013, contains
$r_n=\prod_{i=\lceil n/2\rceil}^{n-1}i$. At $n=2$ that product is $1$,
whereas $r_2=2$. The corrected interval is
\eqref{eq:ratio-product}. The other factorial and ordered-selection
forms in that block agree with $r_n$.
The SageMath expression in the Program field has the same endpoints.
In rising-factorial notation $x^{\overline\ell}:=\prod_{j=0}^{\ell-1}(x+j)$,
its corrected form is
$(\lfloor n/2\rfloor+1)^{\overline{\lceil n/2\rceil}}$.

\section{Coverage and verification}
The two signed source lines say ``Conjecture:''; their dates are
24 November 2012 (A081124) and 30 November 2012 (A084261).
The source versions are A081124 revision~20 (5 November 2025),
A084261 revision~29 (7 September 2020), A081123 revision~62,
and A081125 revision~86. The table distinguishes these conjectures,
previously recorded formulas and additional consequences.
The original posing edits are revisions~8 and~9, respectively;
revisions~9 and~10 approved them on the same dates.
\small
\begin{longtable}{@{}>{\raggedright\arraybackslash}p{0.39\linewidth}>{\raggedright\arraybackslash}p{0.54\linewidth}@{}}
\toprule
Source statement & Coverage\\
\midrule
\endhead
A081124, Mathar, 24 Nov. 2012 & Proved: \eqref{eq:a}.\\
A084261, Mathar, 30 Nov. 2012 & Proved: \eqref{eq:c}.\\
A081123, A081124, A084261, Jovovic, 25 Sept. 2003 & Proved: the recorded even/odd series and
\eqref{eq:egfa}--\eqref{eq:egfc}.\\
A081124 counting interpretation & Prior proof cited: Goyt--Pudwell (2012), Theorem~14.\\
A081124, Sequence Machine (2021) & Proved: the recorded partial sums and shorter recurrence, \eqref{eq:relation} and \eqref{eq:shorter}.\\
A088699, Somos; Barry's thesis & Reproved; prior proof Barry (2009), Section~13.2: array sums \eqref{eq:somos}.\\
A081124/A084261, Nuspl (2023) data & The recorded four-term formulas have all-index proofs in Theorem~\ref{thm:recurrences}.\\
A084261 weighted counting interpretation & Prior proof cited: Fang--Hamaker--Troyka (2020/2022), Proposition~4.13.\\
A161556 row sums and row polynomials & Reproved; prior proof Fang--Hamaker--Troyka (2020/2022): $c_n(t)=\sum_kT(n,k)t^k$.\\
A084261, Hanna; A161556, Barry & Proved: the recorded continued fractions, \eqref{eq:jacobi}; the $t=2$ case is reproved from Barry (2009), Example~64.\\
A137704; A084261, Barry's comment & Proved: the recorded Hankel product, \eqref{eq:hankel}.\\
A081123, Hurt's product comment & Proved: the recorded smaller-part product.\\
A081125, Jovovic, Mathar, Luschny, Hurt, Eldar & Proved: the EGF, recurrence \eqref{eq:ratio}, gamma/swinging forms, larger-part product and sums \eqref{eq:reciprocal}.\\
A081125, Hurt's interval product and SageMath expression & Refuted-and-corrected: $n=2$ gives $1\ne2$; replacement \eqref{eq:ratio-product}.\\
Weighted recurrence \eqref{eq:weighted-recurrence} and partial sums \eqref{eq:weighted-sum} & Not posed in these entries; proved.\\
A004526, floor definition & Used; other interpretations not covered.\\
A081124, A084261, A137704 asymptotics & Not covered.\\
A081123 and A081125 continued fractions & Not covered.\\
\bottomrule
\end{longtable}
\normalsize
The relation \eqref{eq:relation} is the $t=1$ case of
\eqref{eq:weighted-sum}; neither is a conjecture asserted in these entries.

The pinned Lean file
\href{https://github.com/the-omega-institute/trureturing/blob/13f4288c2940bd2aa4020cec789956b9265769f5/D5/S1/Recurrence/MatharA081124.lean}{\texttt{MatharA081124.lean}}
defines the same inclusive binomial sum $a_n$. Its declaration
\nolinkurl{D5.S1.Recurrence.MatharA081124.result} has exactly statement
\eqref{eq:a} for all natural $n\geq3$, interpreted in the integers.
The pinned file
\href{https://github.com/the-omega-institute/trureturing/blob/50af3c7111e85c1112ee9dbc6e721adbb759a26b/D5/S1/Recurrence/MatharA084261.lean}{\texttt{MatharA084261.lean}}
defines $c_n$ by its inclusive even-index binomial sum. Its declaration
\nolinkurl{D5.S1.Recurrence.MatharA084261.result} has exactly statement
\eqref{eq:c} for all natural $n\geq3$, in the integers.
The general transfer lemma, generating functions, weighted recurrence,
array identities, continued fractions, determinants and factorial-ratio
identities have the written proofs above.
The accompanying \texttt{verify.py} checks the initial values and
recurrence-generated terms against their defining sums, the parity
identities, the generating-function coefficients and \eqref{eq:relation}
using exact arithmetic. The accompanying \texttt{verify-siblings.py}
checks the weighted formulas, continued-fraction coefficients, determinant
products and factorial-ratio identities over finite ranges specified in
the README and checker outputs. The written proofs above establish these
identities for all their stated indices; the two pinned Lean declarations
cover \eqref{eq:a} and \eqref{eq:c}, respectively.

\begin{thebibliography}{12}
\raggedright
\bibitem{oeis-a}
OEIS Foundation Inc., \emph{A081124: Binomial transform of
$\lfloor n/2\rfloor!$}, revision 20.
R.~J.~Mathar's Formula entry, 24 November 2012, originally revision~8;
V.~Jovovic's generating function, 25 September 2003.
\url{https://oeis.org/A081124}.

\bibitem{oeis-c}
OEIS Foundation Inc., \emph{A084261: A binomial transform of factorial
numbers}, revision 29.
R.~J.~Mathar's Formula entry, 30 November 2012, originally revision~9;
V.~Jovovic's generating function, 25 September 2003.
\url{https://oeis.org/A084261}.

\bibitem{oeis-b}
OEIS Foundation Inc., \emph{A081123: $a(n)=\lfloor n/2\rfloor!$},
revision 62; and \emph{A004526: Nonnegative integers repeated},
revision 447. \url{https://oeis.org/A081123};
\url{https://oeis.org/A004526}.

\bibitem{gp}
A.~M.~Goyt and L.~K.~Pudwell,
\emph{Avoiding colored partitions of two elements in the pattern sense},
Journal of Integer Sequences \textbf{15} (2012), Article 12.6.2,
Theorem~14, pp.~13--14.
\url{https://cs.uwaterloo.ca/journals/JIS/VOL15/Goyt/goyt4.pdf}.
Also arXiv:1203.3786v3, 13 June 2012.
\url{https://arxiv.org/abs/1203.3786v3}.

\bibitem{fht}
J.~J.~Fang, Z.~Hamaker and J.~M.~Troyka,
\emph{On pattern avoidance in matchings and involutions},
Electronic Journal of Combinatorics \textbf{29}(1) (2022), P1.39.
\url{https://doi.org/10.37236/10155}.
Version used for the locator: arXiv:2009.00079v3, 12 January 2022.
\url{https://arxiv.org/abs/2009.00079v3}.

\bibitem{oeis-triangle}
OEIS Foundation Inc., \emph{A161556}, Paul Barry, 13 June 2009:
triangle formula, row sums and bivariate continued fraction.
\url{https://oeis.org/A161556}.

\bibitem{oeis-hankel}
OEIS Foundation Inc., \emph{A137704}, Paul Barry, 7 February 2008:
Hankel product formula and comments.
\url{https://oeis.org/A137704}.

\bibitem{oeis-ratio}
OEIS Foundation Inc., \emph{A081125: $n!/\lfloor n/2\rfloor!$}, revision~86.
V.~Jovovic, 24 September 2003; R.~J.~Mathar, 26 November 2012;
P.~Luschny, 7 August 2009 and 28 October 2013;
W.~I.~Hurt, 26 January and 6 June 2013; A.~Eldar, 10 March 2022.
\url{https://oeis.org/A081125}.
Swinging-factorial convention: \url{https://oeis.org/A056040}.

\bibitem{sequence-machine}
Sequence Machine, \emph{A081124}, formula records dated
1 November 2021 (partial sums) and 9 December 2021 (second-order recurrence).
\url{https://sequencedb.net/s/A081124}.

\bibitem{nuspl}
P.~Nuspl, \emph{Guessing recurrence sequences in the OEIS}, polynomial
recurrence data, \texttt{ens15.csv} (download links added 7 July 2023),
rows A081124 and A084261.
\url{https://www3.risc.jku.at/people/pnuspl/CSV_OEIS/dfin/ens15.csv}.

\bibitem{somos}
OEIS Foundation Inc., \emph{A088699}, Michael Somos, 8 October 2003:
symmetric array, recurrence and antidiagonal sums.
\url{https://oeis.org/A088699}.

\bibitem{barry-thesis}
P.~Barry, \emph{A Study of Integer Sequences, Riordan Arrays,
Pascal-like Arrays and Hankel Transforms}, Ph.D. thesis,
University College Cork, July 2009.
Section~13.2, pp.~283--284; Example~64, pp.~53--54.
\url{https://research.setu.ie/en/publications/a-study-of-integer-sequencesriordan-arrays-pascal-like-arraysand--4/}.
Author-posted text:
\url{https://www.academia.edu/106522067/}.
\end{thebibliography}

\bigskip
\noindent\small\textcopyright\ 2026 Haobo Ma and Wenlin Zhang.
Licensed under \href{https://creativecommons.org/licenses/by/4.0/}{Creative
Commons Attribution 4.0 International (CC BY 4.0)}.
\end{document}
